\section{Results and Discussion}

\todo{the lower head might be due to leading edge rouding as said by lock, maybe not due to prerotation as the equations suggest.}

\todo{look at H-Q data and see if the cavitation boundary lock describes can be found there as my current cavitation curves look at a different thing.}

\subsection{Experimental Campaign}

The experimental campaign was conducted in three phases: commissioning, electric motor driven pump testing, then integrated turbine and pump testing. The standalone turbine testing was not conducted due to time constraints.

\subsection{Pump: Head--Flow Characteristics}

There are two main curves that are used to evaluate pump performance: the head-flow curve and the efficiency curve. The head-flow curve shows the achievable head at different flow rates and will be discussed first.


% \begin{figure}[H]
%     \centering
%     \begin{subfigure}[b]{0.4\textwidth}
%         \centering
%         
\includegraphics[width=\textwidth]{placeholder.png}
%         \caption{TODO}
%         % \label{fig:TODO}
%     \end{subfigure}
%     \hfill % Adds flexible space between figures
%     \begin{subfigure}[b]{0.4\textwidth}
%         \centering
%         
\includegraphics[width=\textwidth]{placeholder.png}
%         \caption{TODO}
%         % \label{fig:TODO}
%     \end{subfigure}
%     \caption{TODO}
%     % \label{fig:TODO}
% \end{figure}

% \begin{figure}[H]
%     \centering
%     
\includegraphics[width=0.5\textwidth]{placeholder.png}
%     % \label{fig:TODO}
% \end{figure}

% \begin{figure}[H]
%     \centering
%     \begin{subfigure}[b]{0.4\textwidth}
%         \centering
%         
\includegraphics[width=\textwidth]{placeholder.png}
%         \caption{Dimensional head-flow curve}
%         \label{fig:dim-hq}
%     \end{subfigure}
%     \hfill % Adds flexible space between figures
%     \begin{subfigure}[b]{0.4\textwidth}
%         \centering
%         
\includegraphics[width=\textwidth]{placeholder.png}
%         \caption{Non-dimensional head-flow curve}
%         \label{fig:non-dim-hq}
%     \end{subfigure}
%     \caption{TODO}
%     % \label{fig:TODO}
% \end{figure}

\begin{figure}[H]
    \centering
    \includegraphics[width=1\textwidth]{theory_HQ_tuned.pdf}
    \caption{TODO}
    \label{fig:dim-hq} % label moved after caption -- before, \cref bound to the section, not the figure
\end{figure}


\todo{add explanation that RPM was not constant, so affinity laws were used to adjust the data to the average rpm}

% \cref{fig:dim-hq,fig:non-dim-hq}

\cref{fig:dim-hq} Shows the actual achieved head rise across flow rates at the different RPMs that were tested. Overlayed are solid and dashed lines. Both are theoretical head-flow curves generated from \cref{eq:lock_static_head_flow}. The dashed line is with exactly the parameters from Lock's paper, and the solid line is with the pre-rotation $K$ increased from \valLockKOneOriginal to \valLockKOneModified. The dashed vertical dashed line shows the point at a synchronized video camera detected cavitation within the diffuser through the transparent diffuser wall. It can be seen that the original parameters overpredict the head rise, and there are many mechanisms that could be causing this. 

\subsection{Pump: Efficiency and Churning}

The efficiency curve generated by measuring shaft power, subtracting mechanical losses can be shown in the following graph:

The right graph shows the predicted efficiency curve from pure theory, and the right graph shows the curve but with adjusted disc friction. To explain the disc friction adjustment, the pump data at zero flow where all the shaft power was only going into hydraulic losses:

\begin{figure}[H]
    \centering
    \begin{subfigure}[b]{0.49\textwidth}
        \centering
        \includegraphics[width=1\textwidth]{results_eta_flow.pdf}
        \caption{TODO}
        \label{fig:eff} % label moved after caption
    \end{subfigure}
    % \hfill % Adds flexible space between figures
    \begin{subfigure}[b]{0.49\textwidth}
        \centering
        \includegraphics[width=\textwidth]{results_churning_power.pdf}
        \caption{TODO}
        \label{fig:churning}
    \end{subfigure}
    \caption{TODO}
    % \label{fig:TODO}
\end{figure}

\cref{fig:churning} shows that the measured disc friction losses are about \valPumpChurnMult times higher than barske's formula for disc friction. Adjusting the disc friction by this factor leads to a much better fit of the efficiency curve, as can be seen in \cref{fig:eff}.

\begin{figure}[H]
    \centering
    
\includegraphics[width=0.6\textwidth]{placeholder.png} % swap to eta_re_fit.pdf once promoted from plots_round2.py into thesis_figures.py
    \caption{TODO} % eta-Re fit: (1-eta) ~ Re^-a, a = 0.098 +/- 0.010 -> eta(20k) = 26.4% (25.7-27.1%)
    \label{fig:eta-re}
\end{figure}

\subsection{Pump: Cavitation}

Specifically for rocket engine applications, the cavitation behaviour is of high importance. The following graph shows the cavitation curve generated by plotting the head at which cavitation was observed at different flow rates:


\begin{figure}[H]
    \centering
    \includegraphics[width=0.7\textwidth]{results_npshr_vs_q.pdf}
    \caption{TODO}
    \label{fig:npshr}
\end{figure}

   
\cref{fig:npshr} shows the NPSH required to avoid cavitation at different flow rates. The dashed line is the theoretical curve generated by solving \cref{eq:lock_diffuser_throat} for the flow rate at which the head at the diffuser throat is equal to the vapor pressure head. Two interesting observations here.


\todo{add picture of no cavitation, cavitation of throat only, cavitation of inlet, and both}
\todo{3\% vs full cavitation?}

\begin{figure}[H]
    \centering
    \begin{subfigure}[b]{0.45\textwidth}
        \centering
        
\includegraphics[width=\textwidth]{placeholder.png}
        \caption{TODO} % no cavitation
        \label{fig:cav-none}
    \end{subfigure}
    \hfill
    \begin{subfigure}[b]{0.45\textwidth}
        \centering
        
\includegraphics[width=\textwidth]{placeholder.png}
        \caption{TODO} % throat cavitation only
        \label{fig:cav-throat}
    \end{subfigure}

    \begin{subfigure}[b]{0.45\textwidth}
        \centering
        
\includegraphics[width=\textwidth]{placeholder.png}
        \caption{TODO} % eye/inlet cavitation only
        \label{fig:cav-eye}
    \end{subfigure}
    \hfill
    \begin{subfigure}[b]{0.45\textwidth}
        \centering
        
\includegraphics[width=\textwidth]{placeholder.png}
        \caption{TODO} % both sites, descending NPSH
        \label{fig:cav-both}
    \end{subfigure}
    \caption{TODO} % Fig 5.6: annotated cavitation stills -- eye vs throat, descending NPSH (the money shot; from cav_video.py + cav_annotations/)
    \label{fig:cav-stills}
\end{figure}

\subsection{Coupled Turbopump}

% STUB -- ~5,400 rpm settled, measurable pressure rise, deep off-design (p01 ~2.9 bar ~ 38% design);
% NO single forward turbine number (no air mdot meter -> under-determined);
% degeneracy broken by the starting analysis: Mw3 = 1.58 > 1.408 at test speed -> rotor UNSTARTED
% (threshold ~15,300 rpm) -> ~0.25 Nm measured vs ~0.87 Nm started, the 4x gap explained;
% Stark nozzle-attachment check rules out the nozzle; experimental peers block (Donnelly/Peng/Moffitt/Ohlsson).

\begin{figure}[H]
    \centering
    \includegraphics[width=0.7\textwidth]{results_coupled_HQ.pdf}
    \caption{TODO} % Fig 5.7: coupled H-Q scatter vs motor-driven curves
    \label{fig:coupled-hq}
\end{figure}

\begin{figure}[H]
    \centering
    \includegraphics[width=0.7\textwidth]{results_coupled_torque.pdf}
    \caption{TODO} % Fig 5.8: coupled torque vs speed -- data vs started model (dashed) vs unstarted (solid + band)
    \label{fig:coupled-torque}
\end{figure}

\subsection{Extrapolation to the Design Point}

% STUB -- affinity on the pooled psi-phi master curve + Lock at 20k + design star;
% the psi-phi collapse IS the licence for affinity extrapolation -- say so;
% state plainly: coupled test ran ~5,400 rpm, rotor unstarted -> stage similitude NOT matched at design;
% what IS validated: choked nozzle, MoC profile, Goldman starting, affinity/Re-collapsed pump curves;
% two named gaps to 20k: started-stage eta unverified; gamma = 1.4 air vs ~1.2 hot gas.

\begin{figure}[H]
    \centering
    \includegraphics[width=0.7\textwidth]{results_HQ_extrap_20k.pdf}
    \caption{TODO} % Fig 5.9: H-Q affinity extrapolation to 20k + design point
    \label{fig:hq-extrap}
\end{figure}

\subsection{Synthesis}

\begin{table}[H]
  \centering
  \caption{TODO} % predicted vs measured summary -- the quantitative anchor; percentage differences computed by the pipeline from the registry (export_tables.py COMPARE_TABLES, add rows there)
  \label{tab:model_vs_meas}
  % AUTO-GENERATED by export_tables.py (tab_model_vs_meas) — DO NOT EDIT BY HAND.
% Regenerate:  python build_values.py && python export_latex.py && python export_tables.py
\begin{tabular}{@{}lllrl@{}}
\toprule
Quantity & Predicted & Measured & $\Delta$ & Note \\
\midrule
head coefficient $\psi$ & \valPumpPsiDesign & \valPumpPsiMeas & $+3\,\%$ & Barske design vs measured (5-speed collapse) \\
overall efficiency $\eta$ at \valPumpRpmDesign & \valPumpEtaDesign & \valPumpEtaExtrap & $+16\,\%$ & design estimate vs Re-fit extrapolation \\
shutoff churning power & Barske disk correlation & \valPumpChurnMult$\times$ & $+150\,\%$ & measured at shutoff, 15 points \\
rotor-inlet relative Mach $M_{w3}$ at test & \valTurbStartLimit & \valTurbMwTest & $+12\,\%$ & vs Goldman starting limit $\rightarrow$ unstarted \\
coupled shaft torque & \valTurbTorqueStartedTest & \valCoupledTorque & $-69\,\%$ & started-flow model vs measured (unstarted rotor) \\
\bottomrule
\end{tabular}

\end{table}

% STUB (the chapter must build to this) -- toolchain adequacy verdict (Lock within X%/Y pp -> design-stage screening;
% turbine codes validated where speed-independent; stage eta awaits a started test);
% Barske suitability verdict (the central question) -- viable/promising/needs-development, and WHY;
% mechanical redesign learning point (plastic creep -> metal redesign) -- design content, no incident narrative.